\cotefile{91}
\scanpage{87}

\margin{[sous-catégorie pleine de (\unsure{Sch|Schs})]}
\diagram{\scriptsize
\begin{tikzcd}[ampersand replacement=\&,row sep=large]
{\text{Préschémas atomiques : (sommes de schémas de corps)}} \arrow[d, "{\substack{\text{foncteur \unsure{pleinement|pleinmt} fidèle,}\\\text{commutant aux sommes}}}"]\\
{\substack{\text{Foncteurs covariants}\quad(\text{Corps})\unsure{\longrightarrow|\text{---}}(\text{Ens})\\
\text{i.e. contravariants}\quad(\text{Schs ponctuels})^{o}\longrightarrow(\text{Ens})}}
\end{tikzcd}}

\diagram{\scriptsize
\begin{tikzcd}[ampersand replacement=\&,row sep=large]
{\substack{\text{Préschémas atomiques sur } k\quad [\text{i.e. sur }\operatorname{Spec}(k)]\quad
\deleted{(At)}/\operatorname{Spec}k}}
\arrow[d,"{\text{foncteur \unsure{pleinement|pleinmt} fidèle}}"]\\
{\substack{\text{Foncteurs covariants}\quad(\text{Extensions de }k)\longrightarrow(\text{Ens})\\
\text{Foncteurs ''}\qquad(\text{Schs ponctuels}/\operatorname{Spec}k)^{o}\longrightarrow(\text{Ens})}}
\end{tikzcd}}

Soit $X$ un préschéma, il définit un préschéma atomique $|X|=\coprod_{x\in X}\operatorname{Spec}k(x)$, qui \deleted{\unsure{est représenté|est égal}}\inserted{\unsure{correspond|représente}} \unsure{à|:} la restriction du foncteur contravariant défini par $X:(\text{Sch})^{o}\longrightarrow(\text{Ens})$, aux schémas ponctuels.

Bien sur, on obtient un foncteur $X\leadsto |X|$, dont la \uline{restriction aux préschémas réduits} \uline{est} \uline{fidèle}.

\deleted{\unsure{principes|principaux}}

Soient $X,Y$ deux préschémas, on appelle \uline{pseudo-morphisme} ou \uline{application constructible} de $X$ dans $Y$, tt morphisme $f:|\unsure{\overline X|X}|\longrightarrow|\unsure{\overline Y|Y}|$ qui a la propriété suivante : Pour tte partie \deleted{quasi-compact}\inserted{ouverte affine} \deleted{\unsure{donné|affine|de}} $U$ de $X$, il existe une partition de $U$ en des $U_i$ \inserted{\uline{loc. fermés}} \uline{constructibles}, tels que $f\mid |U_i|$ soit un \uline{morphisme} ${U_i}_{\text{réd}}\to Y$.
% archive sheet number: 86
